\documentclass[a4paper,11pt]{article}

\usepackage[utf8]{inputenc}
\usepackage[T1]{fontenc}
\usepackage[ngerman]{babel}
\usepackage[a4paper,margin=2.5cm]{geometry}
\usepackage{lmodern}
\usepackage{enumitem}

\begin{document}
	
	\begin{center}
		{\Large \textbf{Beschreibung der aktuellen Ist-Situation}}\\[0.3cm]
		{\normalsize Lieferbestätigung im aktuellen Dispositionsprozess}
	\end{center}
	
	\vspace{0.8cm}
	
	\section*{1. Ausgangssituation}
	
	Der eigentliche Prozess beginnt erst, nachdem ein Auftrag bereits im System vorhanden ist.
	Die Bestellung erfolgt zuvor über bestehende Kanäle, zum Beispiel über Online-Plattformen wie HeizOel24.
	
	Für die interne Bearbeitung liegen anschließend bereits folgende Informationen vor:
	
	\begin{itemize}[leftmargin=1.2cm]
		\item Kunde
		\item Augtrag
		\item Lieferadresse
		\item Produkt
		\item Bestellmenge
	\end{itemize}
	
	\section*{2. Rolle der Disposition}
	
	Der Disponent erstellt auf dieser Grundlage die vorläufige Tourenplanung.
	
	Dabei werden manuell festgelegt:
	
	\begin{itemize}[leftmargin=1.2cm]
		\item Liefertag
		\item konkretes Lieferfenster
	\end{itemize}
	
	Das Lieferfenster wird aktuell nicht vom Kunden gewählt, sondern zunächst intern vorgegeben.
	
	\section*{3. Telefonische Rückbestätigung}
	
	Spätestens zwei Tage vor dem geplanten Liefertermin erfolgt ein telefonischer Kontakt durch den Disponenten.
	
	Dem Kunden wird ein bereits geplanter Termin genannt.
	
	Ziel ist die Klärung, ob dieser Termin passt.
	
	\section*{4. Ablauf bei Nichterreichbarkeit}
	
	Wenn der Kunde beim ersten Anruf nicht erreichbar ist, erfolgt am selben Tag nach einigen Stunden ein zweiter Anruf.
	
	Falls auch der zweite Versuch erfolglos bleibt:
	
	\begin{itemize}[leftmargin=1.2cm]
		\item wird der Auftrag nicht in die aktuelle Tour aufgenommen
		\item und automatisch auf die nächste Tour verschoben
	\end{itemize}
	
	Falls der Kunde beim zweiten Versuch erreicht wird:
	
	\begin{itemize}[leftmargin=1.2cm]
		\item bleibt der Auftrag in der aktuellen Tour, sofern noch Kapazität vorhanden ist
		\item andernfalls erfolgt ebenfalls eine Verschiebung auf die nächste Tour
	\end{itemize}
	
	\section*{5. Ablauf bei Kundenreaktion}
	
	Wenn der Kunde den vorgeschlagenen Termin bestätigt, gilt dieser als verbindlich.
	
	Der Auftrag bleibt final in der geplanten Tour.
	
	Falls der Kunde zum bestätigten Termin nicht anwesend ist, entsteht eine Vertragsstrafe.
	
	\section*{6. Änderungswünsche des Kunden}
	
	Falls der Kunde den vorgeschlagenen Termin nicht akzeptiert, wird kein neuer Termin sofort vergeben.
	
	Typische Fälle sind:
	
	\begin{itemize}[leftmargin=1.2cm]
		\item Lieferung früher gewünscht
		\item Lieferung später gewünscht
		\item Einschränkungen wie nur vormittags
	\end{itemize}
	
	Diese Informationen werden im DISPO-System gespeichert und bei der nächsten Tourenplanung berücksichtigt.
	
	Eine direkte Neuplanung erfolgt erst später.
	
	\section*{7. Systemseitige Dokumentation}
	
	Alle Rückmeldungen werden im DISPO-System dokumentiert.
	
	Dort existieren Statusinformationen zum Bearbeitungsstand des Auftrags.
	
	\section*{8. Hauptcharakteristik des aktuellen Prozesses}
	
	Der aktuelle Prozess basiert vollständig auf manueller telefonischer Abstimmung durch die Disposition.
	
	Ein automatisierter digitaler Workflow für Rückbestätigung, Statusverarbeitung oder Folgeaktionen existiert derzeit nicht.

\end{document}